\chapter{Conjuntos}
\section{Definición de Conjunto}
Un conjunto es una colección bien definida de objetos distintos, denominados elementos, que pueden ser de cualquier naturaleza.

Para definir un conjunto de manera correcta debemos expresar precisamente que elementos están contenidos por en el.

\begin{ejem}{Conjuntos}
 \begin{itemize}
	\item $A =$ conjunto de todas las camisetas de color rojo.
	\item $B =$ conjunto de los continentes del mundo.
	\item $C =$ conjunto de todos los números pares.
\end{itemize}
\end{ejem}

Los conjuntos se denotan mediante letras mayúsculas ($A, B, C, \dots$), mientras que sus elementos se denotan con letras minúsculas ($a, b, c, \dots$).

\subsection{Pertenencia} 
Afirmar que $x$ es un elemento del conjunto $A$ se formaliza mediante la proposición:
\begin{equation}
	x \in A
\end{equation}
Si el elemento $x$ no pertenece al conjunto $A$, escribimos $x \notin A$. La definición de un conjunto por extensión se realiza enumerando sus elementos entre llaves y cada elemento se separa mediante comas:
\begin{equation}
A = \{a_1, a_2, \dots, a_n\}
\end{equation}
Los conjuntos de pueden definir por extensión o por comprensión.
\subsection{Por extensión}
Cuando definimos un conjunto por extensión se escriben todos sus elementos separados por comas, sin repetirlos y sin importar el orden.
\begin{ejem}{conjuntos por extensión}
	Sean A, B, C y D definidos por extensión.
	\begin{itemize}
		\item $A = \{a,b,c,d,e,f\}$
		\item $B = \{1,2,3,4,5,6,7\}$
		\item $C = \{2,4,6,8,10\}$
		\item $D = \{\{a,b\},\{c,d\},e,f\}$, notar que $\{a,b\}$ y $\{c,d\}$ son elementos del conjunto $D$
	\end{itemize}
\end{ejem}

\subsection{Por comprensión}
Para definir un conjunto por comprensión se deben expresar las características para que un elemento pertenezca a dicho conjunto. Los conjuntos definidos por comprensión tienen la siguiente estructura.
\begin{equation}
A = \{ x \in U \ | \ P(x)\}
\end{equation}
La expresión se lee de la siguiente manera:
\begin{itemize}
	\item $x$ un elemento cualquiera
	\item $\in U$ ese elemento está contemplado en el conjunto $U$
	\item $|$ tal que
	\item $P(x)$ los elementos cumplen las características de la función proposicional $P(x)$
\end{itemize}
Los elementos del conjunto $A$ son los $x$ que pertenecen al conjunto $U$ tal que $P(x)$.

\begin{ejem}{Conjunto expresado por comprensión}
	Dado el conjunto A definido por comprensión:
\begin{itemize}
	\item $A = \{x \in \mathbb{N} \ | \ x < 3 \}$
\end{itemize}
Sabemos que $x$ es un número natural $\mathbb{N}= \{1,2,3,4,5,\dots\}$, veamos que elementos pertenecen al conjunto $A$:
\begin{itemize}
	\item $1 \in \mathbb{N}$, y $1<3$ es $V$, entonces $1 \in A$.
	\item $2 \in \mathbb{N}$, y $2<3$ es $V$, entonces $2 \in A$.
	\item $3 \in \mathbb{N}$, y $3<3$ es $F$, entonces $3 \notin A$.
\end{itemize}
Sabemos que los siguientes números naturales son mayores que $3$ por tanto es absurdo probar con los siguientes números naturales.
\end{ejem}

\begin{dfn}{Conjunto universal}
El conjunto universal es el conjunto más grande dentro de una discusión. A veces el conjunto universo se sobreentiende del contexto o puede ser necesario expresarlo 


\end{dfn}

